\section{Private kernels, signed Palm voids, and the absolute-cumulant obstruction}\label{supp:sec:palm-complements}
This appendix is a methodological complement for the same binary model.
It retains three distinct conclusions: regular configurations drawn from
the target, finite signed Palm identities, and a counterexample to a
global absolute-cumulant smallness condition. All arguments are given
below. The one-sided Palm deficit follows \emph{from}
\cref{thm:micro-run-tv}, never the reverse. None of these constructions
is an input to \cref{supp:sec:microscopic-pivots}.
Keep the hard scales, cutoff $E_*$, maximal support diameter $Q$, original
good set $G_0$ and conditional law $\PA$ used there, and put $\beta=\log2$.

\subsection{Odd rough kernels and a stronger good set}\label{supp:palm:subsec:regular}
Define
\[
 r_Y(m)=\prod_{\substack{q>Y\\v_q(m)\text{ odd}}}q,
 \quad Z_Y(s)=\zeta(2s)\prod_{q\le Y}(1+q^{-s}),
 \quad B(X,Y;s)=X^{s-1}Z_Y(s),\quad\tfrac12<s<1.
\]
The product $r_Y(m)$ and the single pivot $\kappa_{\rm odd}(m)$ have different roles. Every integer has a unique form $m=a^2vr$, with $v$ squarefree and $Y$-smooth and $r=r_Y(m)$ squarefree with primes above $Y$. If
$D_Y(x)=\#\{m\le x:r_Y(m)=1\}$, Rankin gives $D_Y(x)\le x^sZ_Y(s)$.
The exact convolutions over squarefree rough $r$ imply
\begin{align}
 \frac1X\#\{m\le X:r_Y(m)\le T\}
 &\le B(X,Y;s)\left(1+\frac{T^{1-s}}{1-s}\right),\label{supp:palm:eq:kernel-count}\\
 \frac1X\sum_{\substack{m\le X\\r_Y(m)>1}}
       \frac{z^{\omega(r_Y(m))}}{r_Y(m)}
 &\le B(X,Y;s)\left\{\prod_{q>Y}(1+zq^{-1-s})-1\right\}\notag\\
 &\le B(X,Y;s)\{\exp(z/(sY^s))-1\},\qquad z\ge1.\label{supp:palm:eq:kernel-reciprocal}
\end{align}
For the first inequality sum $D_Y(X/r)$ over rough $r\le T$ and enlarge to all positive integers. For the second sum $z^{\omega(r)}D_Y(X/r)/r$; the empty kernel is removed before forming the Euler product. Finally
$\sum_{q>Y}q^{-1-s}\le\sum_{m>Y}m^{-1-s}\le Y^{-s}/s$.

At $s=1-u/V$ and $X=M+O(Q)$, the calculation of Lemma~\ref{supp:pivot:lem:rankin} at $w=V$ gives
\begin{equation}\label{supp:palm:eq:hard-kernel}
 B(X,Y;s)\le e^{-V+o(\nu)},\qquad
 Y^{-s}=e^{-V+u+o(1)}.
\end{equation}
Fix $0<\vartheta<c$, and put
\begin{equation}\label{supp:palm:eq:strong-good}
 T_M=\exp(\vartheta H/u),\qquad
 G_\vartheta=\{j\in G_0:r_Y(j+a)>T_M\ (0\le a\le Q)\}.
\end{equation}
This is a threshold on a product of primes, not a change of the conditioning cutoff $Y$. Eventually $T_M>Y$ and $T_M=M^{o(1)}$.
\begin{proposition}[Cost of stronger deletion]\label{supp:palm:prop:deletion}
One has
\begin{equation}\label{supp:palm:eq:added-deletion}
 p|G_0\setminus G_\vartheta|
 \le \exp\{\log\Lambda-V+\vartheta\nu+o(\nu)\}=o(1),
\end{equation}
and replacing these extra deleted low-type coordinates by independent target coordinates costs at most $2p|G_0\setminus G_\vartheta|$ under $\PA$. Moreover
\begin{equation}\label{supp:palm:eq:cardinality}
 |G_\vartheta^c|\le O(\sqrt M)+M e^{-V+\vartheta\nu+o(\nu)}
                       =o(M/H).
\end{equation}
Completing the original deletions and mark tails has cost
$O(d_{0,M}+p|G_0\setminus G_\vartheta|)=o(1)$.
\end{proposition}
\begin{proof}
Each integer is used in at most $Q+1$ maximal supports.
Use \eqref{supp:palm:eq:kernel-count}--\eqref{supp:palm:eq:hard-kernel} and
$(1-s)\log T_M=\vartheta\nu$; the logarithms of $Q+1$ and $1/(1-s)$ are $O(\log H)=o(\nu)$. This proves the count and the first bound, whose exponent is at most $-I-(c-\vartheta)\nu+o(\nu)$.
At a site of $G_0$, each low-type mean is exactly $p\gamma_\alpha$ on every environment. Thus the added real and target masses are each at most $p$ per site without an extra factor $e^I$.
The original deletion and high marks are covered by \eqref{supp:pivot:eq:deletion}. The last equality in \eqref{supp:palm:eq:cardinality} follows from $V/\nu\sim u\to\infty$ and $\log H=o(\nu)$.
\end{proof}

Call a collection of maximal supports \emph{private-pivot regular} if each raw vertex occurrence has a prime above $Y$, occurring there to an odd exponent, that divides no other vertex occurrence of the collection. Repeated vertices preclude regularity.
\begin{lemma}[CRT allocation of all odd prime factors]\label{supp:palm:lem:crt}
Let $J_1,\ldots,J_k$ be independent uniform indices in $\{1,\ldots,n\}$. Conditional on $J_i=j$ and on a source vertex $m=j+a$ with $r=r_Y(m)>1$, the probability that every prime of $r$ is hosted in one of the other $k-1$ supports is at most
\begin{equation}\label{supp:palm:eq:crt}
 \frac{\{3(k-1)(Q+1)\}^{\omega(r)}}r.
\end{equation}
For $k\ge2$ and $z=3k(Q+1)<Y$, the probability that all indices lie in $G_\vartheta$ but their supports are not regular is at most
$k(Q+1)T_M^{-1+\log z/\log Y}$. A single good support is regular.
\end{lemma}
\begin{proof}
Assign each prime of $r$ to a target block and offset where it is hosted. There are at most $((k-1)(Q+1))^{\omega(r)}$ assignments. For a fixed assignment, CRT fixes each target index modulo the product $r_h$ of its assigned primes. Since $r_h\le r\le M+Q+1<2n$, that residue has probability at most $1/r_h+1/n\le3/r_h$. The target indices are independent and the nontrivial moduli have product $r$. Summing assignments gives \eqref{supp:palm:eq:crt}.
If $r>T_M$, then $\omega(r)\le\log r/\log Y$ and
$z^{\omega(r)}/r\le T_M^{-1+\log z/\log Y}$.
Union-bound over the $k(Q+1)$ source occurrences. A prime above $Y>Q$ cannot be hosted elsewhere inside its own support.
\end{proof}

Let $\mathscr R_M$ be the target configurations of mass at most $K_M=\lceil2\Lambda\rceil$, with low types $e\le E_*$, sites in $G_\vartheta$, and regular maximal supports. The empty configuration is included.
\begin{theorem}[Private pivots for a Poisson-size target cloud]\label{supp:palm:thm:regular}
For the full same-grid target law $\mathsf Q_M$,
\begin{align}\label{supp:palm:eq:regular-error}
 \mathsf Q_M(\mathscr R_M^c)\le\varepsilon_{\rm reg,M}
 \ll{}&\Lambda M^{-1/2}+e^{\log\Lambda-V+\vartheta\nu+o(\nu)}
           +\Lambda2^{-E_*-1}\notag\\
 &+e^{-(2\log2-1)\Lambda}+e^{-c\vartheta\nu^2/(4u)}=o(1).
\end{align}
For every regular $z=\{(j_i,\alpha_i):1\le i\le k\}$,
\begin{equation}\label{supp:palm:eq:presence}
 \PP(\operatorname{Occ}(z)\mid\mathcal F_Y)
       =\prod_{i=1}^k p\gamma_{\alpha_i}
\end{equation}
on every small-prime environment, and therefore under $\PA$.
\end{theorem}
\begin{proof}
Conditional on target mass $N=k$, sites are independent uniform grid indices and types have law $\gamma$. Exponential Markov bounds give
$\PP(N>K_M)\le e^{-(2\log2-1)\Lambda}$.
The bad-site and high-mark expected counts give the first three terms of
\eqref{supp:palm:eq:regular-error}.
For $k\le K_M$, eventually
\[
 1-\frac{\log(3K_M(Q+1))}{\log Y}\ge\frac{c\nu}{2V}.
\]
Apply Lemma~\ref{supp:palm:lem:crt}, obtaining a remaining bound
$K_M(Q+1)e^{-c\vartheta\nu^2/(2u)}$.
Since $\nu^2/(uV)\to\infty$ and $\log(K_M(Q+1))\le V+O(\log H)$, this is at most the final term of \eqref{supp:palm:eq:regular-error}.
For a regular configuration choose one private odd prime per raw vertex. The selected columns of the valuation matrix form an identity minor over $\mathbb F_2$ even conditional on all unchosen signs. The raw values are therefore independent and fair. Each exact signed word fixes $L+e_i+2$ of them; integrating unused maximal-support vertices proves \eqref{supp:palm:eq:presence}.
\end{proof}
The first assertion is about configurations drawn from the target. The second prescribes presence, not absence of further occurrences. For instance, if $X_{2i-1}=X_{2i}=B_i$ with independent Bernoulli-$p$ variables, then target configurations using at most one site per pair have asymptotic probability one when $np^2\to0$ and their prescribed probabilities equal $p^k$. Yet when $np\to\infty$, the actual field lies in that regular set only when empty, of probability $(1-p)^{n/2}\to0$. Its microscopic distance to the Poisson target tends to one. This surrogate is not the arithmetic model.

\subsection{Exact configuration masses and ordinary Palm deletion}\label{supp:palm:subsec:palm}
Put $G=G_\vartheta$, $g=|G|$, $\mu=gp$, and restrict the full arithmetic and target fields to $G$, without truncating exact lengths or filling. Denote their laws by $\mathsf P_M^G,\mathsf Q_M^G$, and their distance by $\Delta_M^G$.
Let $d_A=\PA(\text{a base start occurs outside }G)$ and $\delta_G=p(n-g)$; both tend to zero by Proposition~\ref{supp:palm:prop:deletion} and the original tail estimates. Projection and the deletion coupling give
\begin{equation}\label{supp:palm:eq:projection}
 0\le\Delta_M^{\rm mic}-\Delta_M^G\le d_A+\delta_G.
\end{equation}
Let $\mathscr R_M^G$ be the same regularity conditions inside $G$, with mass bounded by $K_M$. Its target exceptional probability $\varepsilon_{\rm reg,G}=o(1)$ follows by restriction from Theorem~\ref{supp:palm:thm:regular}.
For regular $z$ put $\lambda_z=\prod_{(j,\alpha)\in z}p\gamma_\alpha$ and
\[
 v(z)=\PA(\mathbf I_M=z\mid\operatorname{Occ}(z)),\qquad
 v_G(z)=\PA(\mathbf I_M|_G=z\mid\operatorname{Occ}(z)).
\]
The regular configurations have distinct sites, so
$\mathsf Q_M^G(\{z\})=e^{-\mu}\lambda_z$ and
$\mathsf P_M^G(\{z\})=\lambda_zv_G(z)$.
The normalized retained void is
\begin{equation}\label{supp:palm:eq:normalized-void-definition}
 R_z(t)=\frac{\EA\!\left[
       \prod_{j\in G\setminus\supp z}(1-tJ_{j+1,L})
                       \mid\operatorname{Occ}(z)\right]}
                    {(1-tp)^{g-|z|}},\qquad0\le t\le1.
\end{equation}
It is nonnegative and finite; $p<1$ and the plant is regular.
\begin{proposition}[Averaged deletion and the normalized void deficit]\label{supp:palm:prop:palm}
One has
\begin{align}\label{supp:palm:eq:averaged-deletion}
 &\EE_{\mathsf Q_M^G}\ind_{\mathscr R_M^G}(z)e^\mu(v_G(z)-v(z))\notag\\
 &\hspace{8mm}=\PA(\mathbf I_M|_G\in\mathscr R_M^G,
                     \text{a start occurs outside }G)\le d_A.
\end{align}
For the normalized void $R_z$ in \eqref{supp:palm:eq:normalized-void-definition},
\begin{equation}\label{supp:palm:eq:normalized-deficit}
 \left|\Delta_M^{\rm mic}
 -\EE_{\mathsf Q_M^G}\ind_{\mathscr R_M^G}(z)(1-R_z(1))_+\right|
 \le d_A+\delta_G+\varepsilon_{\rm reg,G}+O(\Lambda^2/n).
\end{equation}
\end{proposition}
\begin{proof}
Presence factorization implies
$\lambda_z(v_G-v)=\PA(\mathbf I_M|_G=z,\mathbf I_M\ne z)$.
Multiplication by the target mass and summation over disjoint events proves
\eqref{supp:palm:eq:averaged-deletion}. For any countable laws,
$\TV(\mathsf P,\mathsf Q)=\sum_z(\mathsf Q(z)-\mathsf P(z))_+$.
Thus, with
$D_G=\EE_{\mathsf Q_M^G}\ind_{\mathscr R_M^G}(1-e^\mu v_G)_+$,
\[
 0\le\Delta_M^{\rm mic}-D_G\le d_A+\delta_G+\varepsilon_{\rm reg,G}.
\]
Once the exact planted words are fixed, no further retained atom is exactly the event that every remaining base-start indicator is zero. Hence
$v_G(z)=(1-p)^{g-k}R_z(1)$, where $k=|z|\le K_M$.
Uniformly on regular configurations,
\[
 |\mu+(g-k)\log(1-p)|
     \le kp+\frac{gp^2}{2(1-p)}=O(\Lambda^2/n).
\]
For $b>0,x\ge0$,
$|(1-bx)_+-(1-x)_+|\le1-e^{-|\log b|}\le|\log b|$;
no bound on $x=R_z(1)$ is needed. This proves \eqref{supp:palm:eq:normalized-deficit}.
\end{proof}
The argument pays $d_A$, not $e^\Lambda d_A$. For the full target, the same mass identity gives directly
\[
 0\le\Delta_M^{\rm mic}
   -\EE_{\mathsf Q_M}\ind_{\mathscr R_M}(z)(1-e^\Lambda v(z))_+
   \le\varepsilon_{\rm reg,M}.
\]
Consequently Theorem~\ref{thm:micro-run-tv} makes both one-sided deficits vanish. In the full-target form, setting
$L_z=(-\log v(z)-\Lambda)_+$ on regular configurations and zero otherwise gives
$(1-e^\Lambda v(z))_+=1-e^{-L_z}$. Its expectation vanishes exactly when $L_z\to0$ in target probability. This is not uniform control of every Palm void, and it does not result from dividing an absolute approximation by $e^{-\Lambda}$.

\begin{corollary}[Completion of the one-sided Palm deficit]\label{supp:palm:cor:palm-completion}
Under Theorem~\ref{thm:micro-run-tv},
\begin{equation}\label{supp:palm:eq:palm-completion}
 \EE_{\mathsf Q^G_M}\left[
        \ind_{\mathscr R^G_M}(z)(1-R_z(1))_+\right]\longrightarrow0.
\end{equation}
Equivalently the lower relative-void defect tends to zero in target-configuration probability. This is an averaged, one-sided statement, not a uniform relative estimate over every plant.
\end{corollary}
\begin{proof}
\Cref{supp:palm:prop:palm} identifies the left side with
$\Delta_M^{\rm mic}+o(1)$, using exact configuration masses and ordinary deletion costs. Apply Theorem~\ref{thm:micro-run-tv}, whose proof does not use that identity.
\end{proof}

\subsection{The finite signed polynomial, including the environment}\label{supp:palm:subsec:signed-void}
Fix a regular plant $z$ and write $\mathcal B_z=G\setminus\operatorname{supp}z$, $m=g-|z|$.
Let $\epsilon$ be the bits of primes at most $Y$, and $\xi$ those above $Y$ among the finitely many queried integers. Successive relative signs give matrices
\[
 J_{j+1,L}=\ind_{\{A_j\epsilon+B_j\xi=b\}},\quad
 b=(1,0,\ldots,0)^{\mathsf T}\in\mathbb F_2^L,
 \qquad A_z\epsilon+B_z\xi=b_z
\]
for the base event and the planted raw equations, respectively. Regularity makes $B_z$ full row rank
$q_z=\sum_{(j,e,\sigma)\in z}(L+e+2)$.
Planting has conditional probability $2^{-q_z}$ on every environment, and so leaves the law of $\epsilon\mid\mathcal A_M$ unchanged. Given that environment, $\xi$ is uniform on the displayed affine solution space.
Put
$\widehat\mu_A(s)=\EE[(-1)^{s\cdot\epsilon}\mid\mathcal A_M]$.
For $S\subseteq\mathcal B_z$, $S\ne\varnothing$, define
\begin{equation}\label{supp:palm:eq:signed-coefficient}
 C_z(S)=
 \sum_{\substack{u_j\in\mathbb F_2^L\setminus\{0\},\ j\in S\\
                  v\in\mathbb F_2^{q_z},\ \sum_j u_jB_j=vB_z}}
 (-1)^{\sum_j u_j\cdot b+v\cdot b_z}
 \widehat\mu_A\left(\sum_j u_jA_j+vA_z\right),
 \qquad C_z(\varnothing)=1.
\end{equation}
For a given collection $(u_j)$ there is at most one $v$, by row independence of $B_z$.
\begin{proposition}[Signed avoidance and a target-sign bound]\label{supp:palm:prop:signed-void}
Writing $\EE_z$ for the planted conditional expectation, for $0\le t\le1$,
\begin{align}
 \EE_z\prod_{j\in S}(J_{j+1,L}-p)&=p^{|S|}C_z(S),\label{supp:palm:eq:centered-plant}\\
 R_z(t)&=1+\sum_{\varnothing\ne S\subseteq\mathcal B_z}
           \left(-\frac{tp}{1-tp}\right)^{|S|}C_z(S).\label{supp:palm:eq:signed-polynomial}
\end{align}
These finite identities include singleton corrections and all overlaps. For fixed unsigned regular geometry $\zeta$, expand the sign dependence as
\[
 R_{\zeta,\boldsymbol\sigma}(1)=\sum_{D\subseteq[k]}a_D(\zeta)
                          \prod_{i\in D}\sigma_i,
 \qquad\mathcal V(\zeta)=\sum_{D\ne\varnothing}a_D(\zeta)^2.
\]
Under the independent fair target signs,
\begin{equation}\label{supp:palm:eq:walsh}
 \EE_{\boldsymbol\sigma}(1-R_{\zeta,\boldsymbol\sigma}(1))_+
 \le\min\left\{1,(1-a_\varnothing(\zeta))_+
                             +\tfrac12\sqrt{\mathcal V(\zeta)}\right\}.
\end{equation}
\end{proposition}
\begin{proof}
The character expansion is
$J_{j+1,L}-p=p\sum_{u_j\ne0}(-1)^{u_j\cdot b}
 (-1)^{(u_jA_j)\cdot\epsilon+(u_jB_j)\cdot\xi}$.
On the planted affine space a character $w\cdot\xi$ has zero mean unless
$w=vB_z$, in which case its mean is
$(-1)^{v\cdot b_z+(vA_z)\cdot\epsilon}$.
Multiplication and environment averaging prove \eqref{supp:palm:eq:centered-plant}.
Expand
$\prod_j(1-tJ_{j+1,L})/(1-tp)^m
 =\prod_j\{1-t(J_{j+1,L}-p)/(1-tp)\}$ to obtain
\eqref{supp:palm:eq:signed-polynomial}.

In planted block $i$, $b_{z,i}=c_i+\tau_i\mathbf1$ with
$\sigma_i=(-1)^{\tau_i}$, where $c_i$ is one at the two borders and zero inside. Thus grouping terms by the block parities of $v$ gives the Walsh expansion. Orthogonality gives mean $a_\varnothing$ and variance $\mathcal V$.
For any integrable $R$ of mean $a$,
\[
 \EE(1-R)_+\le(1-a)_++\EE(a-R)_+
       =(1-a)_++\tfrac12\EE|R-a|
       \le(1-a)_++\tfrac12\sqrt{\Var R}.
\]
Here $R\ge0$, so the left side is also at most one. This proves \eqref{supp:palm:eq:walsh}.
\end{proof}
For the unconditioned model $\widehat\mu_A(s)=\ind_{\{s=0\}}$; for a general admissible event it is not that indicator. Nor are the planted marginal means assumed to equal $p$. The coefficients and their signs must be retained before any norm is taken. Theorem~\ref{thm:micro-run-tv} controls the deficit in \eqref{supp:palm:eq:normalized-deficit}, but does not establish that
$(1-a_\varnothing)_+$ or $\mathcal V$ vanish separately. No such stronger assertion is needed.

Two other exact formulations are useful for separating sufficient conditions from the concluded deficit. If $W_j=J_{j+1,L}-p$ and
\[
 A_z(t)=\sum_{\varnothing\ne S\subseteq\mathcal B_z}
       \left(\frac{t}{1-tp}\right)^{|S|}
         |\operatorname{cum}_z(W_j:j\in S)|,
\]
then the finite moment-partition expansion gives
$|R_z(t)-1|\le e^{A_z(t)}-1$, including singleton Palm corrections.
Its per-site weight at $t=1$ is $(1-p)^{-1}$, not the global weight $2$ discussed below.
Also, with $N_z=\sum_{j\in\mathcal B_z}J_{j+1,L}$ and the planted law tilted by
$(1-t)^{N_z}/\EE_z(1-t)^{N_z}$ for $t<1$,
\[
 -\frac{d}{dt}\log R_z(t)
     =\frac{\EE_{z,t}N_z}{1-t}-\frac{mp}{1-tp}.
\]
Integration to one is an improper identity, with value $+\infty$ when the void is zero. We assert no smallness estimate for either right-hand side.

\subsection{The absolute conditional pair layer}\label{supp:palm:subsec:pairs}
There is a useful pair estimate independent of the microscopic proof. For two good disjoint prescribed raw words of total length $q_1+q_2=q$, write
$A\epsilon+B\xi=b$, $r_0=2^{-q}$,
$\rho_>=q-\operatorname{rank}B$ and
$\rho_{\rm all}=q-\operatorname{rank}[A\ B]$.
If the full system is compatible, a fraction
$a=2^{\rho_{\rm all}-\rho_>}$ of environments is admissible, and on those environments its conditional probability is $r_02^{\rho_>}$, otherwise zero. Thus
\begin{equation}\label{supp:palm:eq:two-ranks}
 \EE_\epsilon|\PP(B\xi=b-A\epsilon\mid\epsilon)-r_0|
 =r_0\{2^{\rho_{\rm all}}-1+2(1-2^{\rho_{\rm all}-\rho_>})\}.
\end{equation}
If the full system is incompatible the left side is $r_0$, and in both cases it is at most
$r_0\{(2^{\rho_{\rm all}}-1)+2\ind_{\{\rho_> >\rho_{\rm all}\}}\}$.
In the incompatible case $\rho_{\rm all}\ge1$.

For a concrete distinction, take $L=3$, $e=f=0$, both signs positive, $Y=11$, and starts $x=92,y=114$. The ten raw vertices are $91,\ldots,95$ and $113,\ldots,117$. Each individual rough row rank is five, their combined rough rank is seven, and their full rank is ten. Thus the unconditioned joint mass is $2^{-10}$, but its conditional value is $2^{-7}$ on one eighth of environments and zero elsewhere. The mean absolute conditional covariance is $7/4096$.

Let $H_2^>(M,Q;Y)$ count ordered separated good maximal-support pairs with nonzero rough nullity. Such a relation uses a vertex $m$ in one support, and every prime of $r_Y(m)$ must be hosted by the other support. Assign those primes to target offsets and use CRT. With $X=M+Q+1$ this gives
\[
 H_2^>(M,Q;Y)\ll(Q+1)X
  \sum_{\substack{m\le X\\r_Y(m)>1}}
              \frac{(Q+1)^{\omega(r_Y(m))}}{r_Y(m)}.
\]
There is only one target block, whose residue count is $O(X/r_Y(m))$ since $r_Y(m)\le X$.
Equations \eqref{supp:palm:eq:kernel-reciprocal}--\eqref{supp:palm:eq:hard-kernel} at $z=Q+1$ imply
\begin{equation}\label{supp:palm:eq:rough-host}
 H_2^>(M,Q;Y)\le M^2e^{-2V+o(\nu)}.
\end{equation}
Indeed $(Q+1)Y^{-s}\to0$ and $u+\log Q=o(\nu)$.
\begin{proposition}[The full signed pair activity]\label{supp:palm:prop:pair-activity}
Uniformly for deterministic $G\subseteq G_0$, the low-type pair activity obeys
\begin{align}\label{supp:palm:eq:pair-activity}
 \mathcal K_{2,M}
 &:=2\sum_{\substack{j<k\\j,k\in G}}\sum_{\alpha,\alpha'\in\mathcal T_*}
          |\Cov_{\mathcal A}(I_{j,\alpha},I_{k,\alpha'})|\notag\\
 &\ll e^Ip^2\{nQ+R^{\rm val}_{2,U}(M,Q+1)+H_2^>(M,Q;Y)\}\notag\\
 &\ll_\varepsilon e^{I+2\log\Lambda-2V+o(\nu)}+M^{-1/3+\varepsilon}=o(1).
\end{align}
\end{proposition}
\begin{proof}
Good-word marginal means are constant on environments. Hence absolute covariance under $\mathcal A_M$ is at most $e^I$ times mean absolute covariance over the unconditioned small-prime environments. Apply \eqref{supp:palm:eq:two-ranks}. Zero extension bounds the full-kernel excess by the maximal full-value profile; the second term requires a rough relation host. Local pairs cost $O(e^Ip^2nQ)$ with signed factor at most four. Sum $\gamma_\alpha\gamma_{\alpha'}$ before applying \eqref{supp:palm:eq:rough-host} and the profile. The budget bounds the first exponential by $e^{-I-2c\nu+o(\nu)}$.
\end{proof}
This is an absolute conditional covariance result, not a graph theorem: zero pair covariance does not imply independence from families of other sites. In particular it cannot replace the exact one-word coupling in Appendix~\ref{supp:sec:microscopic-pivots}.

\subsection{The finite cumulant bound and its arithmetic obstruction}\label{supp:palm:subsec:cumulants}
Let $G\subseteq G_0$ be a deterministic retained set and let $\mathcal T_*=\{0,\ldots,E_M^*\}\times\sgnset$. Define the global distinct-site activity
\begin{equation}\label{supp:palm:eq:absolute-activity}
 \mathcal K_{\rm abs,M}
 =\sum_{\substack{S\subseteq G\\|S|\ge2}}2^{|S|-1}
   \sum_{\boldsymbol\alpha\in\mathcal T_*^S}
       \left|\operatorname{cum}(I_{j,\alpha_j}:j\in S)\right|.
\end{equation}
Here the cumulants are under the stated arithmetic law, not under a product target.


For a finite family of categories $Y_j\in\{\partial\}\cup\mathcal T_j$ with marginals $\nu_j$, let $I_{j,a}=\ind_{\{Y_j=a\}}$ and let $\mathcal K$ be the activity \eqref{supp:palm:eq:absolute-activity} with these types and their actual law.
\begin{lemma}[A finite categorical cumulant bound]\label{supp:palm:lem:cumulant}
One has
\begin{equation}\label{supp:palm:eq:cumulant-tv}
 \TV\left(\Law((Y_j)_j),\bigotimes_j\nu_j\right)
             \le\min\{1,(e^{2\mathcal K}-1)/2\}.
\end{equation}
For $t\ge0$, if $B_j$ are Bernoulli variables of means $r_j$, and
$F_t=\EE\prod_j\{1+t(B_j-r_j)\}$, then
\begin{equation}\label{supp:palm:eq:polynomial-envelope}
 |F_t|\le\exp\!\left\{
   \sum_{\substack{S\\|S|\ge2}}t^{|S|}
                |\operatorname{cum}(B_j:j\in S)|\right\}.
\end{equation}
For categorical occupancies $B_j=\sum_a I_{j,a}$ with $2\max_jr_j<1$,
$\mathcal K\ge\tfrac12\log F_2$.
\end{lemma}
\begin{proof}
The exact identity
$\delta_{Y_j}=\nu_j+\sum_a(I_{j,a}-\EE I_{j,a})(\delta_a-\delta_\partial)$
gives, after tensor expansion, centered moment coefficients on subsets $S$.
The signed tensor has zero-mass supremum norm $2^{|S|-1}$. Singletons vanish.
Expand each centered moment into cumulant partitions, with blocks of size at least two. If
$b_B=2^{|B|-1}\sum_{\boldsymbol a}|\operatorname{cum}(I_{j,a_j}:j\in B)|$,
a partition into $m$ blocks contributes at most
$2^{m-1}\prod_B b_B$. Sum disjoint block collections, then drop disjointness to obtain
$\frac12\sum_{m\ge1}(2\sum_B b_B)^m/m!=(e^{2\mathcal K}-1)/2$.
The identical finite partition expansion with weights $t^{|B|}$ proves
\eqref{supp:palm:eq:polynomial-envelope}. Multilinearity gives
$2\mathcal K\ge\sum_{|S|\ge2}2^{|S|}|\operatorname{cum}(B_j:j\in S)|$.
When $2\max r_j<1$, all factors in $F_2$ are positive, so taking logarithms proves the last assertion. No analytic infinite cumulant series is assumed.
\end{proof}
The law on the right of \eqref{supp:palm:eq:cumulant-tv} is the product of categorical marginals, not yet independent Poisson coordinates. Poissonizing occupied mass $s_j\le p$ costs at most $\sum_js_j^2\le np^2$ by coupling the total count and using the same type kernel. The finite inequality remains true, but its global-smallness hypothesis fails as follows.

\begin{theorem}[Small microscopic distance with divergent absolute activity]\label{supp:palm:thm:cumulant-obstruction}
Let $q$ tend to infinity through odd primes and set
\[
 a_q=\lfloor\log_2q\rfloor,\qquad M_q=2^{q-1+a_q},\qquad L_q=q-1.
\]
Take $\mathcal A_{M_q}$ to be the whole space. If the retained sets have the exact low-type marginals and omit $o(M_q/\log M_q)$ sites, then
\begin{equation}\label{supp:palm:eq:cumulant-obstruction}
 \liminf_{q\to\infty}\frac{\log M_q}{M_q}\mathcal K_{\rm abs,M_q}
 \ge\frac{\log2}{2}(\log3-1)>0.
\end{equation}
This applies to $G_0$ of Appendix~\ref{supp:sec:microscopic-pivots} and to the stronger retained sets in Appendix~\ref{supp:sec:palm-complements}. The two-site contribution tends to zero, so the activity of orders at least three diverges. Nevertheless $\Delta_{M_q}^{\rm mic}\to0$, already within the earlier sufficient labelled range.
\end{theorem}
\begin{proof}[Proof of Theorem~\ref{supp:palm:thm:cumulant-obstruction}]
For $M=M_q$, $L=q-1$, $a_q=\lfloor\log_2q\rfloor$,
\[
 H\sim\beta q,\qquad
 \Lambda=2^{a_q}-(q-1)2^{-(q-1)}\asymp q.
\]
Thus $a_q=o(\log M)$ and both $\log\Lambda$ and
$\log\Lambda+\log(1+\Lambda)$ are $O(\log H)=o(V-c\nu)$ for every fixed $c$.
Let $g=|G|$, $S_*=1-2^{-E_*-1}$ and
$B_j=\sum_{\alpha\in\mathcal T_*}I_{j,\alpha}$.
Their common exact mean is $r=pS_*\le p$, and $gr\le\Lambda$.

Consider the event
\[
 \mathcal E_q=\{f(q)=-1,\ f(\ell)=1
              \text{ for every prime }\ell\le M+Q,\ \ell\ne q\}.
\]
Its probability is $2^{-\pi(M+Q)}$ in the original unconditioned model; it is not a new admissible conditioning event. On this event every queried value is $(-1)^{v_q(m)}$.
For $(t+1)q\le M$ and $q\nmid t(t+1)$, the endpoints $tq,(t+1)q$ are negative and the $q-1$ interior values are positive. Therefore the site $j=tq$ carries the low type $(0,+1)$.
With $T=\lfloor M/q\rfloor-1$, the number of certified indices is exactly
\[
 T-\lfloor T/q\rfloor-\lfloor(T+1)/q\rfloor
       =(1-2/q)M/q+O(1).
\]
Omitting $o(M/H)=o(M/q)$ sites leaves at least
$N_0=(1-o(1))M/q$ retained low-type atoms.
For $N=\sum_{j\in G}B_j$ and sufficiently large $q$,
\[
 F_2=(1-2r)^g\EE\left(\frac{3-2r}{1-2r}\right)^N
 \ge 2^{-\pi(M+Q)}(1-2r)^g
                   \left(\frac{3-2r}{1-2r}\right)^{N_0}.
\]
The prime number theorem gives $\pi(M+Q)=(1+o(1))M/H$;
$g\log(1-2r)\ge-O(\Lambda)$; and
$\log((3-2r)/(1-2r))=\log3+O(r)$.
Since $\Lambda\asymp H=o(M/H)$, it follows that
\[
 \log F_2\ge\{\beta(\log3-1)+o(1)\}M/H.
\]
Lemma~\ref{supp:palm:lem:cumulant} proves \eqref{supp:palm:eq:cumulant-obstruction}.
The omission condition holds for $G_\vartheta$ by \eqref{supp:palm:eq:cardinality} and for $G_0$ by its smaller deletion. Proposition~\ref{supp:palm:prop:pair-activity} gives the vanishing pair layer, and hence divergence of the remaining orders. These parameters lie even in the graph-based labelled budget. In particular \cref{thm:micro-run-tv} gives the asserted vanishing microscopic distance, independently of the activity calculation.
\end{proof}
For a fixed $t>0$, the same witness has leading lower exponent
$\beta\{\log(1+t)-1\}M/H$. It diverges when $t>e-1$; this is a threshold for the witness, not an optimal boundary of the model. Restricting the occupancy to $N\le b\Lambda$, $b>1$, can have vanishing ordinary total-variation cost, but it changes its means, cumulants, and conditional forcing identities. No such nonlocal truncation is used to retain an invalid absolute-smallness claim.
